% !TEX TS-program = xelatex
% !TEX encoding = UTF-8 Unicode
% !Mode:: "TeX:UTF-8"
% Tailored for: Biological / Environmental Sciences Consultant roles

\documentclass{my-resume}
\usepackage{my-resume-zh}
\usepackage{linespacing_fix}
\usepackage{cite}

\begin{document}
\pagenumbering{gobble}

\photoheader{%
  \name{吴筱琦}
  \contactInfo{+86 13776020628}{yukiwxq@gmail.com}{现居城市：伦敦/苏州}{意向城市：上海/苏州/杭州 \textbar\ 生物/环境科学咨询}
}{../template/photo.jpg}{2.4cm}

\section{教育背景}
\entry{帝国理工学院 (QS 2)}{2025.09 - 2026.09}{计算生命科学研究型硕士}{伦敦}
\entry{帝国理工学院 (QS 2)}{2022.09 - 2025.07}{生物科学 本科}{伦敦}
\begin{itemize}[parsep=0.2ex]
  \item \textbf{相关课程}：生物多样性与保护生物学，全球变化生物学，生态学与进化论，行为生态学，生态学实地技能，生态学基础，非洲生物学野外课程，进化与多样性
\end{itemize}

\section{实习与社会实践}
\entry{上海宣泰医药科技有限公司}{2024.07 - 2024.09}{分析工程师 \textperiodcentered\ 分析部}{上海}
\begin{itemize}[parsep=0.2ex]
  \item \textbf{质量决策支持}：使用 HPLC/UV 对客户提供的药品样本开展稳定性检测，通过严格质控流程系统鉴别杂质风险，为客户提供可直接用于稳定性评估决策的分析数据
  \item \textbf{流程与沟通}：独立完成样本前处理、仪器操作与数据记录全流程，确保检测结果的准确性与可追溯性
\end{itemize}
\entry{上海剑藤教育科技有限公司}{2023.07 - 2023.08}{教学助教 \textperiodcentered\ 生化部}{上海}
\begin{itemize}[parsep=0.2ex]
  \item \textbf{教学执行}：为冲刺牛津、剑桥等名校的学生提供课外拓展教学辅导，独立承担部分全英文授课、模拟面试与课堂讨论
  \item \textbf{需求分析与成果}：整理并分析学生需求文档，为授课教师提供差异化教学建议，最终帮助多名学生拿到剑桥本科录取通知
\end{itemize}

\section{科研与项目经历}

\entry{昆虫衰退的驱动因素、影响与保护对策综述}{2024.03 - 2025.06}{}{伦敦}
\role{文献综述、批判性分析}{}
\begin{itemize}[parsep=0.2ex]
  \item \textbf{证据综合}：系统审阅 70+ 篇文献，分析栖息地丧失、农药使用、气候变化等驱动因素，其中一项全球荟萃分析显示昆虫生物量每年下降约 2.5\%
  \item \textbf{影响评估}：评估昆虫衰退对生态系统服务（传粉、养分循环）、农业与人类健康的影响，形成具政策参考价值的综述报告
  \item \textbf{批判分析}：批判性评估文献中的矛盾证据、"有得有失"反叙事及地理/分类学报告偏差，提出未来研究与政策建议方向
\end{itemize}

\entry{北大西洋渔业数据变化驱动因素分析与食物网构建}{2025.01 - 2025.02}{}{伦敦}
\role{Excel、文献综述、食物网分析}{}
\begin{itemize}[parsep=0.2ex]
  \item \textbf{数据分析}：量化 8 种鱼类 90 年长期渔获数据，识别鲽形目鱼类等因兼捕与晚熟因素面临的种群脆弱性（1990-2000 年间下降约 40\%）
  \item \textbf{建议输出}：结合政策/气候/渔业活动文献综合分析，提出普通海鸦作为指示物种，为渔业资源管理提供科学依据——报告获高分与优秀评价
\end{itemize}

\entry{环境花色比例对昆虫颜色选择的影响}{2024.09 - 2024.10}{}{开普敦}
\role{野外实验、R（t 检验、Pearson 相关）}{}
\begin{itemize}[parsep=0.2ex]
  \item \textbf{实验设计}：在南非野外课程中，于 6 个样地设计五色陷阱盘实验，研究昆虫访花颜色偏好与环境花色比例的关系
  \item \textbf{实验结果}：黄色、蓝色和白色陷阱均显著高于基线吸引昆虫（p$<$0.001）；黄色捕获率与环境黄花比例显著正相关（r=0.369，p=0.027）
\end{itemize}

\section{技能/其他}
\entrySkillsSep
\begin{itemize}[parsep=0.2ex]
  \item \textbf{生态与野外研究}：野外实验设计（如多样地彩色陷阱盘偏好实验）、具观察者间信度检验的行为学测定、生物多样性/环境调查、利益相关方沟通与汇报、宏观生态学与系统发育比较分析（AVONET 级性状数据、PGLS）
  \item \textbf{科研写作与分析}：文献综合与批判性分析、政策相关建议撰写、R/Python 统计支持（GLM、PGLS、Wilcoxon、t 检验）
  \item \textbf{语言}：英语（雅思 7.5），普通话（母语），法语（A2）；\textbf{兴趣爱好}：钢琴，摄影，滑雪
\end{itemize}

\end{document}
